% Artículo 2 — atributos de liberación de los lotes y reactogenicidad de VA-MENGOC-BC (español).
\documentclass[11pt]{article}
\newcommand{\DocLang}{es}
\input{papers/shared/preamble.tex}
\input{papers/shared/authors.tex}
\input{papers/shared/pending.tex}
\input{papers/p2-lot-quality/shared/flow.tex}
\externaldocument{supplement}
\usepackage{lineno}
\onehalfspacing

\hypersetup{pdftitle={Atributos de liberacion de los lotes y reactogenicidad de la vacuna VA-MENGOC-BC},
  pdfauthor={Richard Alejandro Matos Arderi; Abel Ponce Gonzalez}}

\begin{document}

\begin{titlepage}
\SyntheticNotice
\begin{center}
{\LARGE\bfseries Atributos de liberación de los lotes y reactogenicidad de una vacuna de vesículas
de membrana externa: vinculación de los datos de control de calidad de \PTNLots{} lotes de
VA-MENGOC-BC con las notificaciones nacionales de eventos adversos\par}
\bigskip
\AuthorBlock
\end{center}
\bigskip
\noindent\textbf{Título abreviado:} Calidad de los lotes y reactogenicidad de VA-MENGOC-BC\par
\noindent\textbf{Palabras clave:} liberación de lotes de vacunas; control de calidad; control
estadístico de procesos; vesículas de membrana externa; endotoxina; eventos supuestamente
atribuibles a la vacunación o inmunización; diseño de solo casos; farmacovigilancia\par
\end{titlepage}

\linenumbers

\section*{Resumen}
\noindent\textbf{Antecedentes:} Se debate si la calidad medida de los lotes de vacuna influye en las
reacciones notificadas, y no se ha estudiado en las vacunas de vesículas de membrana externa (VME),
cuya reactogenicidad depende de la endotoxina y del adyuvante.\par
\noindent\textbf{Métodos:} Los datos de liberación de \PTNLots{} lotes de VA-MENGOC-BC
(\PTLotFirstYear--\PTLotLastYear; \PTNAttributes{} atributos) se analizaron con control de procesos,
capacidad ($P_{pk}$), monitorización multivariante robusta, detección de puntos de cambio y pruebas
de equivalencia entre lotes nacionales y de exportación. Las notificaciones cubanas de eventos
supuestamente atribuibles a la vacunación o inmunización (\VStudyFirstYear--\VStudyLastYear) se
vincularon con los lotes administrados. En un diseño de solo casos, ecuaciones de estimación
generalizadas agrupadas por lote estimaron si las notificaciones de fiebre o de reacciones locales
severas variaban con los atributos del lote, con control de la tasa de falsos descubrimientos, una
exposición de control negativo y análisis de sensibilidad.\par
\noindent\textbf{Resultados:} \PTConformitySentence; el $P_{pk}$ fue menor que 1,0 en
\PTNAttrPpkBelowOne{} atributos y los lotes nacionales y de exportación fueron equivalentes en
\PTNEquivalent{} de \PTNAttributes. De \PTNTargetReports{} notificaciones, \PTNLinked{}
(\PTLinkRate) se vincularon con \PTNLinkedLots{} lotes. \PTNSignificantFDR{} de \PTNPrimaryTests{}
asociaciones preespecificadas superaron la corrección (\PTSignificantList); la razón de odds por
desviación estándar de endotoxina para la fiebre de al menos 39\,\textdegree{}C fue
\PTEvFeverThreeNineEndotoxinOR{} (\CI{} \PTEvFeverThreeNineEndotoxinCI). \PTNegControlSentence, y
los atributos de los lotes \PTIncrementalWord{} la predicción fuera del lote.\par
\noindent\textbf{Conclusiones:} Vincular los ensayos rutinarios de liberación con la
farmacovigilancia es factible y permite contrastar hipótesis sobre la reactogenicidad relacionada con
los lotes dentro de especificación, en apoyo de la gestión de riesgos para la calidad y de la
vigilancia poscomercialización de las vacunas de VME.

\clearpage
\section{Introducción}\label{sec:intro}

Cada lote de una vacuna se libera solo después de comprobar sus atributos de calidad frente a las
especificaciones de su registro \cite{WHO2013LotRelease}. Dentro de esos límites, atributos como el
contenido de antígeno, la adsorción al adyuvante o la endotoxina residual varían de un lote a otro.
Si esta variabilidad dentro de especificación influye en las reacciones de los vacunados es una
pregunta que los sistemas de calidad farmacéutica \cite{ICHQ9R1,ICHQ10} y la farmacovigilancia rara
vez responden juntos. El interés por la cuestión creció cuando análisis de sistemas nacionales de
notificación sugirieron heterogeneidad de los eventos adversos sospechados entre lotes de vacunas
contra la COVID-19 \cite{Schmeling2023,Manniche2024,Furst2024,Schmeling2025}. Esos análisis mostraron
también los riesgos del enfoque: el número de dosis realmente administradas de cada lote, las
diferencias entre las poblaciones vacunadas y los cambios de la notificación en el tiempo pueden
crear diferencias aparentes entre lotes \cite{Hviid2023,Scott2023}. Trabajos metodológicos han
propuesto estadísticos de barrido sobre la genealogía de fabricación para localizar estas señales
\cite{Mahaux2019}, pero ningún estudio ha relacionado los atributos de liberación medidos con las
reacciones notificadas.

Las vacunas de vesículas de membrana externa (VME) contra la enfermedad meningocócica por
serogrupo~B son un escenario idóneo para esta pregunta. Su reactogenicidad ---sobre todo fiebre y
reacciones locales en niños pequeños--- está bien descrita \cite{Nokleby2007,Holst2009} y tiene
determinantes plausibles a nivel de lote: la endotoxina unida a la membrana de las vesículas
desencadena la respuesta febril \cite{Sheerin2019,Zariri2016}, y la adsorción en hidróxido de
aluminio reduce la pirogenicidad \cite{Rosenqvist1998}. VA-MENGOC-BC, la vacuna cubana de VME contra
los serogrupos~B y~C, se produce de forma continua desde hace más de tres décadas y se administra a
todos los lactantes cubanos \cite{Sotolongo2007,SierraGonzalez2019}. Su fabricante conserva el
registro de liberación de cada lote, y el sistema nacional de vigilancia registra el lote
administrado en cada notificación de ESAVI \cite{Galindo2012}.

Vinculamos estas dos fuentes con dos objetivos: (i) caracterizar la consistencia de la producción de
VA-MENGOC-BC mediante control estadístico de procesos, capacidad del proceso y monitorización
multivariante, y (ii) estimar, en un diseño de solo casos, si los atributos de los lotes se asocian
con el perfil de reacciones notificadas tras la vacunación, con hipótesis preespecificadas, una
exposición de control negativo y una evaluación explícita de la potencia estadística.

\section{Métodos}\label{sec:methods}

\subsection{Diseño y fuentes de datos}\label{sec:design}
El estudio combinó un análisis retrospectivo de los datos de liberación de los lotes con un análisis
de solo casos de las notificaciones de ESAVI vinculadas, y se informa según RECORD-PE
\cite{Langan2018}. El registro de liberación contenía \PTNLots{} lotes producidos en
\PTLotFirstYear--\PTLotLastYear, \PTNNationalLots{} para el programa nacional y \PTNExportLots{} para
exportación, con \PTNAttributes{} atributos de calidad y sus especificaciones. Las notificaciones de
ESAVI procedieron de la base nacional de vigilancia descrita en un análisis complementario:
\VNFiles{} ficheros anuales de \VStudyFirstYear--\VStudyLastYear, armonizados, seudonimizados y
depurados de duplicados con el mismo flujo de análisis de código abierto (versión \VCodeVersion)
\cite{Kahn2016,ENISA2019}. \VFormSentence. Los datos de los lotes son información industrial confidencial; se
analizaron en una estación de trabajo institucional y solo se publican como posición dentro de la
ventana de especificación (0 = límite inferior, 1 = límite superior), nunca como valores brutos.

\subsection{Atributos y escalas}\label{sec:scales}
Los atributos se analizaron en su escala natural o tras la transformación que hace su variación
aproximadamente simétrica: logaritmo en base 2 para el título bactericida, logaritmo en base 10 para
la concentración de IgG anti-VME y $\log_{10}(1+x)$ para la endotoxina. En los análisis de asociación,
cada atributo se estandarizó con la media y la desviación estándar (DE) de los lotes vinculados, de
modo que las razones de odds se refieren a una diferencia de una DE entre lotes.

\subsection{Desempeño del proceso}\label{sec:m-spc}
Para cada atributo, en orden de producción, se construyeron gráficos de valores individuales y de
rangos móviles con las reglas de Nelson \cite{Nelson1984,Montgomery2019} y un gráfico de media móvil
ponderada exponencialmente ($\lambda=\PTEWMALambda$) \cite{Lucas1990}. El desempeño del proceso se
resumió con el $P_{pk}$, con intervalos bootstrap de percentiles a partir de \PTCapBootstrap{}
remuestreos \cite{Efron1987}, y con el $P_{pk}$ basado en percentiles de la norma ISO~22514-2, que no
supone normalidad \cite{ISO22514-2}, para todos los lotes y para los producidos antes
(\PTPeriodOne) y durante (\PTPeriodTwo) la ventana de farmacovigilancia; la división la fijó el
diseño del estudio, no los datos. Los desplazamientos bruscos se localizaron con un algoritmo PELT
exacto con penalización de tipo BIC \cite{Killick2012,Truong2020}. La atipicidad multivariante se
monitorizó con un modelo robusto de componentes principales (determinante de covarianza mínimo
\cite{Rousseeuw1999}, componentes que explican al menos el 85\,\% de la varianza), con límites de
control de la $T^2$ de Hotelling a partir de la distribución beta \cite{Tracy1992} y el error
cuadrático de predicción con límites de Jackson--Mudholkar \cite{Jackson1979,MacGregor1995}, ambos con
$\alpha=\PTMSPCAlpha$. Los lotes nacionales y de exportación se compararon con dos pruebas
unilaterales \cite{Schuirmann1987} (errores estándar de Welch; margen de equivalencia del
\PTEqMarginPct{} de la ventana de especificación), el desplazamiento de Hodges--Lehmann
\cite{Hodges1963} y una prueba de energía multivariante \cite{Szekely2013}.

\subsection{Vinculación de las notificaciones con los lotes}\label{sec:m-link}
Los números de lote registrados en las notificaciones se normalizaron y se emparejaron con el
registro de liberación en tres pasos: coincidencia exacta, coincidencia en el núcleo del código (sin
separadores ni sufijos) y coincidencia aproximada (similitud de al menos \PTFuzzyThreshold{} entre
lotes nacionales producidos antes de la vacunación). Las coincidencias ambiguas se resolvieron a
favor de los lotes nacionales y del lote más reciente producido antes del año de vacunación. Los
enlaces que implicaban una vacunación más de \PTMaxYearsAfterProduction{} años después de la
producción se consideraron inverosímiles y se excluyeron del análisis principal
(\figref{fig:flow}). Se compararon las características de las notificaciones vinculadas y no
vinculadas para evaluar el sesgo de vinculación (\tabref{tab:S-linkbias} del material
suplementario).

\begin{figure}[tbp]
\centering
\LinkageFlow
\caption{Vinculación de las notificaciones de ESAVI con VA-MENGOC-BC con el registro de liberación
de los lotes administrados. Las cifras proceden de la publicación con control de divulgación.}
\label{fig:flow}
\end{figure}

\subsection{Análisis de asociación de solo casos}\label{sec:m-assoc}
Como no se registra el número de dosis administradas de cada lote, no pueden estimarse riesgos
absolutos por lote. Por ello se utilizó un diseño de solo casos \cite{Piegorsch1994}: entre las
notificaciones vinculadas con un lote, se estimó si la probabilidad de que una notificación incluyera
una reacción determinada variaba con los atributos del lote. El diseño es válido para el perfil
relativo de las reacciones bajo el supuesto de que la exhaustividad de la notificación de una
reacción respecto a otra no varía con los atributos del lote.

La familia principal preespecificada comprendió tres desenlaces (fiebre de al menos 39 y
40\,\textdegree{}C y reacción local severa) y cuatro exposiciones con un mecanismo plausible
(endotoxina, hidróxido de aluminio, concentración de proteínas y adsorción de proteínas), es decir,
\PTNPrimaryTests{} asociaciones. El volumen de llenado, que no debería influir en el tipo de
reacción, fue una exposición de control negativo \cite{Lipsitch2010,Arnold2016}. Cada asociación se
estimó con una ecuación de estimación generalizada logística con correlación de trabajo
intercambiable dentro del lote y errores estándar robustos \cite{Liang1986,Zeger1986}, ajustada por
edad (logaritmo de meses), sexo, número de dosis, coadministración y año calendario; también se
presentan las estimaciones sin ajustar. El procedimiento
de Benjamini--Hochberg controló la tasa de falsos descubrimientos en \PTFDRAlpha{} dentro de la
familia principal \cite{Benjamini1995}; el control negativo no formó parte de la familia. Los demás
atributos, la atipicidad multivariante ($T^2$) y los desenlaces secundarios fueron exploratorios.

Los análisis de sensibilidad restringieron los datos a los enlaces exactos, excluyeron las vacunas
coadministradas y conservaron los enlaces temporalmente inverosímiles; reajustaron las exposiciones en
un modelo logístico bayesiano con intercepto aleatorio, y modelaron a nivel de lote la proporción de
notificaciones con cada reacción con un modelo cuasibinomial ponderado por el número de
notificaciones por lote. Para interpretar los resultados nulos, se calculó la razón de odds mínima
detectable por DE con una potencia del 80\,\% a partir del número de notificaciones, la prevalencia
de cada reacción y la correlación intralote (efecto de diseño) \cite{Hsieh1998}.

\subsection{Valor predictivo incremental}\label{sec:m-ml}
Se evaluó si los atributos de un lote ayudan a predecir las reacciones notificadas con él más allá
de las características del vacunado. Se entrenaron árboles con potenciación del gradiente
\cite{Ke2017} con y sin los atributos de los lotes y se evaluaron por el área bajo la curva ROC en
una validación cruzada de \PTCVFolds{} particiones agrupadas por lote, de modo que cada predicción
se refiere a lotes no vistos durante el entrenamiento. La diferencia de áreas se comparó con su
distribución bajo \PTNPermutations{} permutaciones de los perfiles completos de atributos entre lotes,
lo que conserva la correlación entre atributos y la agrupación de las notificaciones.

\subsection{Control de divulgación y ética}\label{sec:m-sdc}
Solo salieron del entorno controlado resultados agregados: se suprimieron los recuentos no nulos
menores que \VMinCell{} junto con los estadísticos derivados de ellos, los atributos de los lotes se
publicaron solo como posición dentro de la ventana de especificación y un verificador automático
revisó cada fichero publicado \cite{Hundepool2012}. El estudio utilizó registros de vigilancia
seudonimizados y registros de fabricación según el procedimiento institucional de gobierno de datos
del Instituto Finlay de Vacunas y la Ley 149/2022 \cite{Ley149}.

\section{Resultados}\label{sec:results}

\subsection{Consistencia de la producción}\label{sec:r-spc}
\PTConformitySentence{} (menor porcentaje de lotes conformes en cualquier atributo,
\PTPctConforming). El $P_{pk}$ fue de al menos 1,33 en \PTNAttrPpkAbove{} atributos y menor que 1,0
en \PTNAttrPpkBelowOne{} (\PTPpkBelowOneList; \tabref{tab:capability}, \figref{fig:capability});
los valores menores que 1,0 indican que el centrado o la dispersión del proceso hacen esperables
lotes cercanos a un límite. El algoritmo PELT localizó \PTNChangepoints{} desplazamientos sostenidos
(\PTChangepointList). El modelo robusto de componentes principales retuvo \PTNComponents{}
componentes que explican el \PTExplainedVar{} de la varianza; \PTTTwoFlagged{} lotes superaron el
límite de la $T^2$ y \PTSPEFlagged{} el del error cuadrático de predicción (\figref{fig:t2}). Los
lotes nacionales y de exportación fueron equivalentes dentro del margen preespecificado en
\PTNEquivalent{} de \PTNAttributes{} atributos (\tabref{tab:S-equivalence} del material
suplementario), y sus distribuciones multivariantes \PTEnergyWord{} (prueba de energía
$p$\PTEnergyP).

\begin{table}[tbp]
\centering\small
\caption{Desempeño del proceso de los lotes de VA-MENGOC-BC: $P_{pk}$ con intervalo bootstrap del
95\,\% para todos los lotes y por período de producción, y año de producción de los puntos de cambio
localizados por PELT.}
\label{tab:capability}
\PTTableCapability
\end{table}

\begin{figure}[tbp]
\centering
\includegraphics[width=\textwidth]{p2_qc_annual}
\caption{Atributos de liberación de los lotes de VA-MENGOC-BC por año de producción, expresados como
posición dentro de la ventana de especificación (0, límite inferior; 1, límite superior). Las líneas
muestran la mediana anual y la banda el rango intercuartílico; las líneas rojas marcan los límites
de especificación.}
\label{fig:qcannual}
\end{figure}

\begin{figure}[tbp]
\centering
\begin{minipage}[t]{0.49\textwidth}\centering
\includegraphics[width=\linewidth]{p2_capability}
\caption{$P_{pk}$ (intervalo bootstrap del 95\,\%) por período de producción. Las líneas verticales
marcan 1,00 y 1,33.}
\label{fig:capability}
\end{minipage}\hfill
\begin{minipage}[t]{0.49\textwidth}\centering
\includegraphics[width=\linewidth]{p2_t2}
\caption{Atipicidad multivariante ($T^2$) de los lotes por año de producción, con el límite de
control.}
\label{fig:t2}
\end{minipage}
\end{figure}

\subsection{Vinculación}\label{sec:r-link}
De \PTNTargetReports{} notificaciones con VA-MENGOC-BC, \PTNLinked{} (\PTLinkRate) se vincularon con
\PTNLinkedLots{} lotes: \PTLinkExact{} por coincidencia exacta, \PTLinkCore{} por el núcleo del
código y \PTLinkFuzzy{} por coincidencia aproximada (\figref{fig:flow}; \tabref{tab:S-linkage} del
material suplementario). El análisis principal incluyó \PTNAnalysed{} notificaciones en
\PTNAnalysedLots{} lotes (mediana de \PTReportsPerLot{} notificaciones por lote).

\subsection{Atributos de los lotes y reacciones notificadas}\label{sec:r-assoc}
\PTNSignificantFDR{} de las \PTNPrimaryTests{} asociaciones preespecificadas siguieron siendo
significativas tras la corrección por la tasa de falsos descubrimientos (\PTSignificantList;
\tabref{tab:gee}, \figref{fig:forest}). La proporción de notificaciones con fiebre de al menos
39\,\textdegree{}C \PTEvFeverThreeNineEndotoxinDir{} el contenido de endotoxina del lote
(razón de odds por DE \PTEvFeverThreeNineEndotoxinOR, \CI{} \PTEvFeverThreeNineEndotoxinCI;
$q$\PTEvFeverThreeNineEndotoxinQ) y \PTEvFeverThreeNineAlohThreeConcDir{} su contenido de hidróxido de
aluminio (\PTEvFeverThreeNineAlohThreeConcOR, \CI{} \PTEvFeverThreeNineAlohThreeConcCI). La fiebre de
al menos 40\,\textdegree{}C \PTEvFeverFourZeroEndotoxinDir{} la endotoxina
(\PTEvFeverFourZeroEndotoxinOR, \CI{} \PTEvFeverFourZeroEndotoxinCI), y la proporción de
reacciones locales severas \PTEvSevereLocalReactionAlohThreeConcDir{} el hidróxido de aluminio
(\PTEvSevereLocalReactionAlohThreeConcOR, \CI{} \PTEvSevereLocalReactionAlohThreeConcCI). \PTNegControlSentence{} ($p<0{,}05$ en \PTNNegControlNominal{} de tres desenlaces). \PTSensSentence{} (\tabref{tab:S-sensitivity} del
material suplementario).

\begin{table}[tbp]
\centering\small
\caption{Asociación entre los atributos de los lotes y las reacciones notificadas con VA-MENGOC-BC
(diseño de solo casos). Razones de odds por DE del atributo entre lotes, de ecuaciones de estimación
generalizadas logísticas agrupadas por lote, crudas y ajustadas por edad, sexo, dosis,
coadministración y año; $q$, valor de $p$ ajustado por Benjamini--Hochberg dentro de la familia preespecificada (el
control negativo queda fuera de la familia).}
\label{tab:gee}
\PTTableGEE
\end{table}

\begin{figure}[tbp]
\centering
\includegraphics[width=\textwidth]{p2_forest}
\caption{Razones de odds ajustadas por DE de cada atributo de los lotes (IC 95\,\%) para cada
reacción, modelos de exposición única. La banda sombreada abarca el efecto mínimo detectable (de 1/EMD
a EMD); los marcadores grises muestran la exposición de control negativo.}
\label{fig:forest}
\end{figure}

\subsection{Potencia y valor predictivo incremental}\label{sec:r-power}
Con el número de lotes y la correlación intralote observados, la razón de odds mínima detectable por
DE fue \PTEvFeverThreeNineMDE{} para la fiebre de al menos 39\,\textdegree{}C,
\PTEvFeverFourZeroMDE{} para la fiebre de al menos 40\,\textdegree{}C y
\PTEvSevereLocalReactionMDE{} para las reacciones locales severas (\tabref{tab:power}). Añadir todos
los atributos de los lotes a las características del vacunado \PTIncrementalWord{} la predicción
fuera del lote de las reacciones: el cambio del área bajo la curva ROC fue
\PTEvFeverThreeNineDeltaAUC{} para la fiebre de al menos 39\,\textdegree{}C ($p$ de permutación
\PTEvFeverThreeNinePermP; \figref{fig:incremental}).

\begin{table}[tbp]
\centering\small
\caption{Efecto mínimo detectable (EMD; razón de odds por DE con potencia del 80\,\%, $\alpha=0{,}05$
bilateral) dada la agrupación de las notificaciones en lotes, y valor incremental de los atributos de
los lotes para la predicción fuera del lote ($\Delta$AUC, cambio del área bajo la curva ROC; $p$ de
las permutaciones de los perfiles de atributos entre lotes).}
\label{tab:power}
\PTTablePower
\end{table}

\begin{figure}[tbp]
\centering
\includegraphics{p2_incremental}
\caption{Cambio del área bajo la curva ROC con validación cruzada al añadir los atributos de los
lotes (puntos), frente a su distribución de permutación (se marca el percentil 95).}
\label{fig:incremental}
\end{figure}

\section{Discusión}\label{sec:discussion}

\subsection{Hallazgos principales}\label{sec:d-principal}
El registro de liberación de \PTNLots{} lotes de VA-MENGOC-BC mostró un proceso
\PTWithinSpecPhrase, con \PTNChangepoints{} desplazamientos sostenidos y una capacidad menor que 1,0
en \PTNAttrPpkBelowOne{} atributos. Al vincular el \PTLinkRate{} de las notificaciones de ESAVI con
sus lotes, \PTNSignificantFDR{} de \PTNPrimaryTests{} asociaciones preespecificadas entre los
atributos de los lotes y el perfil de reacciones notificadas superaron la corrección por
multiplicidad (\PTSignificantList), mientras que la exposición de control negativo
\PTNegControlWord{} con las reacciones.

\subsection{Interpretación}\label{sec:d-interpretation}
Una asociación entre endotoxina y fiebre sería biológicamente coherente: el lipooligosacárido
asociado a las VME activa la inmunidad innata a través del receptor de tipo Toll~4, y la respuesta
febril a las vacunas de VME se le ha atribuido \cite{Sheerin2019,Zariri2016}, mientras que la
adsorción en hidróxido de aluminio modula la pirogenicidad \cite{Rosenqvist1998}. Toda asociación de
este tipo debe leerse según lo que el diseño puede mostrar: un cambio en la composición de las
notificaciones entre lotes que estaban todos dentro de especificación, no un riesgo absoluto por
dosis. La exposición de control negativo, los análisis de sensibilidad y los efectos mínimos
detectables ayudan a separar los efectos reales de los lotes de los artefactos que han afectado a los
análisis por lote de notificaciones espontáneas: diferencias entre las poblaciones vacunadas con cada
lote, en el tiempo calendario y en la notificación \cite{Hviid2023,Scott2023}. Por la misma razón, la
ausencia de asociación solo es informativa para efectos mayores que las razones de odds mínimas
detectables aquí presentadas.

El control estadístico de procesos aplicado retrospectivamente a los datos de liberación aporta una
evidencia distinta: si el proceso se mantuvo estable, si cambió bruscamente y si los lotes para uso
nacional y para exportación fueron equivalentes. Un índice de capacidad menor que 1,0 no implica no
conformidad, pero indica que el margen entre el proceso y sus límites es estrecho y es un objetivo
natural de la gestión de riesgos para la calidad \cite{ICHQ9R1,ICHQ10}. Combinar estos indicadores con
la vinculación de seguridad convierte el registro de liberación de los lotes en una herramienta de
vigilancia poscomercialización.

\subsection{Fortalezas y limitaciones}\label{sec:d-limitations}
El estudio vincula por primera vez, en una vacuna de VME, la calidad medida de cada lote con las
reacciones notificadas tras su administración; utiliza el registro de liberación completo y las
notificaciones nacionales de \VStudyFirstYear--\VStudyLastYear, preespecifica sus hipótesis y
controla la multiplicidad, y genera cada número a partir de una publicación verificada y con control
de divulgación.

Tiene limitaciones importantes. En primer lugar, se desconocían las dosis administradas por lote,
por lo que el diseño estima el perfil relativo de las reacciones y no su incidencia; un atributo que
aumentara por igual la frecuencia de todas las reacciones no se detectaría. En segundo lugar, la
exhaustividad de la notificación podría diferir entre lotes de formas relacionadas con sus
atributos, por ejemplo a través del tiempo calendario; el ajuste por año y el control negativo
atenúan, pero no excluyen, este sesgo. En tercer lugar, en algunas notificaciones el número de lote
faltaba o no se encontró, y las notificaciones vinculadas y no vinculadas pueden diferir
(\tabref{tab:S-linkbias} del material suplementario). En cuarto lugar, los lotes estuvieron
\PTWithinSpecPhrase, por lo que los resultados no se extrapolan fuera del intervalo especificado. En
quinto lugar, los ensayos de liberación miden una muestra de cada lote en el momento de la
liberación y no recogen los cambios durante el almacenamiento y la distribución. Por último, las
reacciones se registraron como indicadores del formulario nacional, sin validación clínica.

\subsection{Implicaciones}\label{sec:d-implications}
Los fabricantes ya custodian los registros de liberación y los sistemas nacionales ya registran el
lote administrado; el flujo de análisis utilizado aquí muestra que ambos pueden vincularse de forma
rutinaria, preservando la confidencialidad de los datos industriales mediante el control de
divulgación. Registrar las dosis distribuidas por lote permitiría estimar la incidencia por lote, y
el mismo marco puede aplicarse a otras vacunas con datos de calidad por lote.

\section{Conclusiones}\label{sec:conclusions}
Los lotes de VA-MENGOC-BC se produjeron \PTWithinSpecPhrase. Vincular los atributos de liberación de
los lotes con las notificaciones nacionales de ESAVI es factible y ofrece un modo reproducible de
contrastar si la variación dentro de especificación se asocia con la reactogenicidad; las
asociaciones identificadas aquí (\PTSignificantList) definen hipótesis para su confirmación y para la
gestión de riesgos para la calidad.

\section*{Declaraciones}
\input{papers/p2-lot-quality/shared/declarations_es.tex}

\bibliographystyle{vancouver}
\bibliography{papers/shared/bib/pubmed,papers/shared/bib/manual}

\end{document}
